\subsection{Test System and Data}
The water operator runs the SNET network of the repository: one ground-level reservoir, one fixed-speed pump
(design point 50~L/s at 40~m, wire-to-water efficiency 0.75, calibrated on-state power $\overline{P}_p=25.95$~kW),
two junctions, two pipes, and one elevated cylindrical tank (12~m diameter, 25~m
elevation, level 0--6~m, initial 3~m). Daily
demand is 2,214~m$^3$ with the diurnal pattern of Fig.~\ref{fig:inputs}; a 20~m pressure
floor and terminal replenishment are enforced. The tank holds about 2.8 pump-hours of water, so
the shiftable volume is tank-limited in the sense of the two-period model. The energy operator owns the two-bus feeder of the
repository with its daily load shape scaled to a 100~kW peak, a 200~kW$_\mathrm{p}$
PV array, a 200~kWh/50~kW battery ($\eta=0.95$, initial and minimum terminal energy 100~kWh), and a connection rated
150~kW import and 50~kW export. Day-ahead purchase prices are 0.08 in hours 0, 1, 2, 3, 4, 5, 6, 21, 22, and 23; 0.16 in hours 7, 8, 9, 10, 11, 12, 13, 14, and 15; 0.32 in hours 16, 17, 18, 19, and 20
(unit/kWh); the sale price is 0.03. The flat retail tariff of the baseline is $\bar\tau=0.16$, the
time average of the purchase price. Profiles are synthetic and illustrative (Table~\ref{tab:case}).

\begin{table}[!t]
\caption{Case-Study Data (Synthetic Profiles, One-Hour Intervals)}
\label{tab:case}
\centering
\footnotesize
\setlength{\tabcolsep}{4pt}
\begin{tabular}{@{}lp{5.4cm}@{}}
\toprule
Item & Value \\
\midrule
Water network & SNET: 1 reservoir, 1 pump, 2 junctions, 2 pipes, 1 tank \\
Tank & 12~m diameter, elevation 25~m, level 0--6~m, initial 3~m \\
Demand, pressure floor & 2,214~m$^3$/day with daily pattern; 20~m \\
Pump & 50~L/s at 40~m; $\eta_p=0.75$; $\overline{P}_p=25.95$~kW; 6 PWL segments \\
Pipes, line currents & 12 PWL segments (signed domain), 5 segments \\
Feeder & Two-bus circuit, load shape peaking at 100~kW \\
PV, battery & 200~kW$_\mathrm{p}$; 200~kWh, 50~kW, $\eta=0.95$, $e^0=100$~kWh \\
Connection & 150~kW import, 50~kW export; $0.9\le v\le1.1$~p.u. \\
Prices (unit/kWh) & purchase 0.08/0.16/0.32 (off/mid/on-peak); sale 0.03; flat tariff 0.16 \\
\bottomrule
\end{tabular}
\end{table}

\begin{figure}[!t]
\centering
\includegraphics[width=\columnwidth]{inputs.pdf}
\caption{Input profiles: day-ahead purchase and sale prices, water demand, feeder load, and available PV power.}
\label{fig:inputs}
\end{figure}

\subsection{Experiments}
Six experiments are run: (E1) the flat-tariff baseline and the two tariff benchmarks against the integrated optimum;
(E2) the payoffs of the revenue-sharing contract, the fee window, and a numerical deviation test; (E3) the replay of the
baseline and coordinated schedules in EPANET and OpenDSS; (E4) the price coordination scheme with the bundle
master started from the energy operator's marginal costs and from the market price, and with the subgradient rule of the two-agent model; (E5) sweeps of the import and export
ratings of the connection; and (E6) the restricted-LP internal prices that verify Proposition~\ref{prop:price}.
All models were solved with HiGHS 1.11.0 through Pyomo; the integrated MILP has 3219
variables and 2598 constraints and was solved to a 0.003\% relative gap in
76~s, the coordination subproblems to a 0.1\% gap. Timing includes the
Pyomo interface overhead and is not a solver benchmark.

\subsection{The Value of Coordination (E1, E6)}
Table~\ref{tab:arrangements} compares the arrangements. Under the flat tariff the water operator needs
9 pump-hours and, holding its level, places 2~h at 0.08, 5~h at 0.16, 2~h at 0.32 (unit/kWh);
the site's day-ahead cost is 107.94/day. The integrated optimum pumps 3~h at 0.08, 6~h at 0.16, 0~h at 0.32 and costs 90.91/day,
a coordination gain $\mathcal{G}=17.03$/day, 15.8\% of the site's bill and
46\% of the water operator's baseline pumping bill of 37.37/day.
Figure~\ref{fig:computed} shows how the gain is earned: the coordinated schedule places 5 of its 9 pump-hours in the export and idle hours, where the site's own energy is cheapest, and none in the peak block at 0.32; the baseline, holding its level, places 4 in the surplus hours and 2 in the peak block, which the grid serves at the peak price.

The internal price $\pi^\star_t$ of the restricted LP follows Proposition~\ref{prop:price}: in the 15 hours in which the site imports it equals the purchase price; in the 6 export hours (hours 9, 10, 11, 12, 13, and 14) it equals the sale price 0.03; in the 3 idle hours (hours 7, 8, and 16) it lies inside the band $[c^{\mathrm{sell}}_t, c^{\mathrm{buy}}_t]$, set by the battery's intertemporal arbitrage. Passing the purchase price through to the water operator captures 95.6\% of the gain: the water operator then avoids the peak block but places only 5 pump-hours in the surplus hours, having no reason to prefer them over the other hours of the same purchase price. Charging the internal price itself captures 98.2\%: the water operator is indifferent among equally priced hours, and the tie it happened to break is not the one the energy operator can serve most cheaply (Proposition~\ref{prop:tariff}).

\begin{table}[!t]
\caption{Arrangements Compared (Unit/Day). Gain and Capture Are Relative to the Baseline}
\label{tab:arrangements}
\centering
\scriptsize
\setlength{\tabcolsep}{2pt}
\begin{tabular}{@{}lrrrrrr@{}}
\toprule
Arrangement & Cost & Gain & Capt.\ (\%) & On (h) & Import & Export \\
\midrule
Baseline: flat tariff, level holding & 107.94 & 0.00 & 0.0 & 9 & 856 & 212 \\
Pass-through tariff $\tau_t=c^{\mathrm{buy}}_t$ & 91.66 & 16.28 & 95.6 & 9 & 830 & 201 \\
Internal-price tariff $\tau_t=\pi^\star_t$ & 91.21 & 16.73 & 98.2 & 9 & 830 & 216 \\
Integrated optimum (contract) & 90.91 & 17.03 & 100.0 & 9 & 830 & 226 \\
\bottomrule
\end{tabular}
\end{table}

\begin{figure}[!t]
\centering
\includegraphics[width=\columnwidth]{coordination_results.pdf}
\caption{Baseline versus coordinated operation. Top to bottom: purchase and sale prices with the internal price
$\pi^\star_t$; pump load; tank level (markers are inventory boundaries, the dotted line the initial level); net import
at the connection (negative is export).}
\label{fig:computed}
\end{figure}

\subsection{Dividing the Gain (E2)}
Table~\ref{tab:payoffs} lists the payoffs. Because the site is a net importer, the settlement is negative
($R^\star=-90.91$/day), and the contract is read as bill sharing: with the fee anchored on the baseline,
$F(\varphi)=\mathrm{bill}^0+\varphi R^0$, the water operator's payoff rises by $\varphi\mathcal{G}$ and the energy
operator's by $(1-\varphi)\mathcal{G}$ at every share (Corollary~\ref{cor:window}). At the equal split $\varphi=1/2$
the fee is -16.60/day and each operator gains 8.51/day; for that share any fee in
$[-25.11,\,-8.09]$ leaves both operators no worse off than the baseline, an interval of width
$\mathcal{G}$. The deviation test in Table~\ref{tab:deviation} evaluates the water operator's contract payoff at
$\varphi=1/2$ if it ran one of the other schedules instead, with the energy operator responding optimally: every
deviation loses, as Theorem~\ref{thm:contract} requires.

\begin{table}[!t]
\caption{Daily Payoffs by Arrangement (Unit/Day); Fee Anchored on the Baseline}
\label{tab:payoffs}
\centering
\scriptsize
\setlength{\tabcolsep}{2.5pt}
\begin{tabular}{@{}lrrrrrrr@{}}
\toprule
Arrangement & $\varphi$ & $F$ & $\Pi_{\mathrm{W}}$ & $\Pi_{\mathrm{E}}$ & $\Pi$ & $\Delta\Pi_{\mathrm{W}}$ & $\Delta\Pi_{\mathrm{E}}$ \\
\midrule
Baseline (flat tariff) & -- & -- & -37.37 & -70.57 & -107.94 & -- & -- \\
Contract & 0.00 & 37.37 & -37.37 & -53.54 & -90.91 & +0.00 & +17.03 \\
Contract & 0.50 & -16.60 & -28.85 & -62.05 & -90.91 & +8.51 & +8.51 \\
Contract & 1.00 & -70.57 & -20.34 & -70.57 & -90.91 & +17.03 & +0.00 \\
\bottomrule
\end{tabular}
\end{table}

\begin{table}[!t]
\caption{Water Operator's Contract Payoff at $\varphi=1/2$ for Alternative Schedules (Unit/Day)}
\label{tab:deviation}
\centering
\footnotesize
\setlength{\tabcolsep}{4pt}
\begin{tabular}{@{}lrr@{}}
\toprule
Schedule run by W & $\Pi_{\mathrm{W}}$ & Deviation gain \\
\midrule
Baseline: flat tariff, level holding & -37.37 & -8.51 \\
Pass-through tariff $\tau_t=c^{\mathrm{buy}}_t$ & -29.23 & -0.38 \\
Internal-price tariff $\tau_t=\pi^\star_t$ & -29.00 & -0.15 \\
Integrated optimum (contract) & -28.85 & +0.00 \\
\bottomrule
\end{tabular}
\end{table}

\subsection{Distributed Computation (E4)}
Table~\ref{tab:coordination} and Fig.~\ref{fig:convergence} summarize the price coordination runs. Started from the
energy operator's marginal cost of serving the metered baseline schedule, the first message of the recommended protocol, the
scheme certifies a gap of 0.72\% after the first exchange, the recovered schedule attains the
integrated optimum at iteration 2, and the bundle master stops after 7 iterations, when its model
predicts no further improvement at the exchanged price resolution, with a dual bound of 90.55/day and a
certified gap of 0.39\%. Started from the market price instead, the first exchange certifies only
8.18\%, because the purchase price leaves the energy operator indifferent in every import hour;
the scheme attains the optimum at iteration 7, reaches a 1\% certificate at iteration 17,
and ends after 29 iterations with a gap of 0.02\%. The subgradient rule from the
marginal-cost start attains the optimum at iteration 2 and reaches a gap of 0.39\% after
15 iterations. Evaluating the dual function at the restricted-LP prices $\pi^\star$ gives $D(\pi^\star)=90.91$/day against $C^\star=90.91$/day: the Lagrangian relaxation of the interface has no duality gap on this instance (Proposition~\ref{prop:convergence}(iv)), so the residual certified gap of the runs is the price of finite iterations, quantized prices, and subproblem tolerances, not a structural limit. For comparison, the central solver's own bound on the integrated MILP is
90.90, a gap of 0.003\%, and it holds both operators' data. Each iteration exchanges 48 scalars; no demand forecast, tank state, pump
characteristic, PV forecast, battery state, or market position crosses the boundary.

\begin{table}[!t]
\caption{Price Coordination Runs: Iterations, Messages, Certified Dual and Primal Values, Final Gap, the Iteration at Which the Recovered Schedule First Attains the Integrated Optimum, and Wall Time}
\label{tab:coordination}
\centering
\scriptsize
\setlength{\tabcolsep}{2pt}
\begin{tabular}{@{}lrrrrrrr@{}}
\toprule
Run & Iter. & Msgs & Dual & Primal & Gap \% & Opt.\ at & Min \\
\midrule
Bundle, marginal-cost start & 7 & 336 & 90.55 & 90.91 & 0.39 & 2 & 4.5 \\
Bundle, market-price start & 29 & 1392 & 90.89 & 90.91 & 0.02 & 7 & 57.2 \\
Subgradient, marginal-cost start & 15 & 720 & 90.55 & 90.91 & 0.39 & 2 & 10.7 \\
\bottomrule
\end{tabular}
\end{table}

\begin{figure*}[!t]
\centering
\includegraphics[width=\textwidth]{convergence.pdf}
\caption{Left: certified gap between the recovered primal value and the certified dual bound per iteration for the bundle
master from the marginal-cost and market-price starts and for the subgradient rule. Right: internal-price iterates at three
hours for the bundle run from the market-price start; thin lines mark the restricted-LP prices $\pi^\star_t$.}
\label{fig:convergence}
\end{figure*}

\subsection{Connection-Limit Sensitivities (E5)}
Table~\ref{tab:sweep} and Fig.~\ref{fig:sensitivity} vary the ratings of the connection. Tightening the import rating lowers the gain from 17.03/day at 150~kW to 10.92/day at 70~kW; at 100~kW and below the internal price departs from the market price in hours 21, 22, and 23, where the rating binds. No schedule is feasible at 60~kW. Tightening the export rating lowers the gain from 17.03/day at 50~kW to 16.61/day at 0~kW: with less room to export, both schedules absorb more of the midday surplus on site, and the internal price in the curtailment hours falls to the marginal value of spilled energy, zero, for the coordinated schedule and the baseline alike.

\begin{table}[!t]
\caption{Connection-Limit Sweeps (Unit/Day). ``Off Reference'' Lists the Non-Idle Hours in Which the Internal Price Departs from the Market Price of the Connection's Direction}
\label{tab:sweep}
\centering
\footnotesize
\setlength{\tabcolsep}{3pt}
\begin{tabular}{@{}lrrrrp{2.2cm}@{}}
\toprule
Limit & kW & Baseline & Integrated & Gain & Off reference \\
\midrule
Import & 150 & 107.94 & 90.91 & 17.03 & none \\
Import & 125 & 107.94 & 90.91 & 17.03 & none \\
Import & 100 & 110.28 & 98.67 & 11.61 & hours 21, 22, and 23 \\
Import & 90 & 116.54 & 105.12 & 11.43 & hours 21, 22, and 23 \\
Import & 80 & 122.99 & 111.75 & 11.24 & hours 5, 21, 22, and 23 \\
Import & 70 & 129.57 & 118.64 & 10.92 & hours 5, 6, 21, 22, and 23 \\
Import & 60 & infeasible & infeasible & -- & -- \\
Export & 50 & 107.94 & 90.91 & 17.03 & none \\
Export & 25 & 110.26 & 93.18 & 17.09 & hours 9, 10, 11, 12, 13, and 14 \\
Export & 0 & 114.29 & 97.68 & 16.61 & none \\
\bottomrule
\end{tabular}
\end{table}

\begin{figure*}[!t]
\centering
\includegraphics[width=\textwidth]{sensitivity.pdf}
\caption{Coordination gain versus the import rating (left) and the export rating (right) of the connection. Shaded bands
mark ratings at which only coordinated operation is feasible or no schedule is feasible.}
\label{fig:sensitivity}
\end{figure*}

\subsection{Independent Nonlinear Replay (E3)}
The baseline and coordinated pump schedules are replayed in EPANET with a five-minute hydraulic time step and hourly
reports, the level controls of the input file being replaced by the hourly commands, and the optimized net bus injections
are solved as balanced constant-power snapshots in OpenDSS. Table~\ref{tab:replay} summarizes the discrepancies.
Both schedules satisfy the pressure floor and the voltage band at every report time and the imposed pump statuses match the
MILP decisions. The replay ends with the tank 0.137~m (baseline) and 0.124~m (coordinated) below its initial level despite the replenishment constraint \eqref{eq:w_tank_terminal}, so the hourly hydraulic approximation does not preserve the terminal requirement exactly. The electrical discrepancies are small because the feeder is short and lightly loaded,
the regime in which the flat-voltage linearization is most accurate. Report-time checks do not certify feasibility between
reports; the results support the schedules as planning-stage outputs, not as a certification for field operation.

\begin{table}[!t]
\caption{Replay of the Baseline and Coordinated Schedules in EPANET and OpenDSS}
\label{tab:replay}
\centering
\footnotesize
\setlength{\tabcolsep}{4pt}
\begin{tabular}{@{}lrr@{}}
\toprule
Metric & Baseline & Coordinated \\
\midrule
Minimum junction pressure (m) & 25.92 & 25.03 \\
Pressure violations at report times & 0 & 0 \\
Pump-status mismatches at report times & 0 & 0 \\
Maximum tank-head discrepancy (m) & 0.139 & 0.125 \\
Maximum flow discrepancy (L/s) & 0.350 & 0.374 \\
Terminal tank-head change vs.\ initial (m) & $-0.137$ & $-0.124$ \\
All OpenDSS snapshots converged & yes & yes \\
Maximum voltage discrepancy (p.u.) & 0.00002 & 0.00003 \\
Voltage violations & 0 & 0 \\
Peak line loading (\%) & 0.17 & 0.19 \\
Maximum grid-exchange discrepancy (kW) & 0.000 & 0.000 \\
\bottomrule
\end{tabular}
\end{table}
